\documentclass{article}

\setlength{\textwidth}{6.5in}
\setlength{\textheight}{9.0in}
\setlength{\oddsidemargin}{-.525cm}
\setlength{\evensidemargin}{-.525cm}
%\addtolength{\topmargin}{-1.0in}
\setlength{\parindent}{.7cm}

\begin{document}

%{Biographical Sketches}
\def\namefont{\Large\bf}
\def\itemfont{\Large\sc}

\newdimen\leftdim
\leftdim=108pt

\newdimen\rightdim
\rightdim=324pt

\newdimen\wholedim
\wholedim=432pt

\newdimen\middim
\middim=12pt

\newdimen\smalldim
\smalldim=2pt

\begin{center}
{\Large\sc Hongwei Xi} \\[6pt]
Computer Science Department \\
Boston University \\
Boston, MA~02215 \\
Email: {\tt hwxi@cs.bu.edu} \\
Tel: +1 (617) 358-2511 \kern12pt Fax: +1 (617) 353-6457 \\
URL: {\tt http://www.cs.bu.edu/\char126{hwxi}} \\
\end{center}

\begin{tabular}{rp{\rightdim}}
\hbox to \leftdim{\itemfont Education\hss} & \\[\middim]
{December 1998} & {\underline{PhD} in Pure \& Applied Logic, Carnegie Mellon University} \\
                & {Dissertation: {\em Dependent Types in Practical Programming}} \\
                & {Thesis Advisor: {Frank Pfenning}} \\[\smalldim]

{July 1988}     & {\underline{MS} in Mathematics, Nanjing University} \\[\smalldim]
{June 1985}     & {\underline{BS} in Mathematics, Nanjing University} \\[\middim]

\hbox to \leftdim{\itemfont Professional Experience\hss} & \\[\middim]
Sept. 2007 & {\underline{Associate Professor}, Computer Science Department, Boston University} \\[\smalldim]
2001 -- 2007 & {\underline{Assistant Professor}, Computer Science Department, Boston University} \\[\smalldim]
1999 -- 2001 & {\underline{Assistant Professor}, Department of Electrical \& Computer Engineering and Computer Science, University of Cincinnati} \\[\smalldim]
1998 -- 1999 & {\underline{Post-doctoral Research Associate}, Department of Computer Science and Engineering, Oregon Graduate Institute} \\[\smalldim]
1992 -- 1998 & {\underline{Teaching and Research Assistant}, Carnegie Mellon University} \\[\smalldim]
1988 -- 1992 & {\underline{Assistant Professor}, Department of Computer Science and Engineering, Shanghai Jiao Tong University, China} \\[\middim]

\hbox to \leftdim{\itemfont Project Related Publications\hss} & \\[\middim]

\multicolumn{2}{c}{%
\begin{minipage}{\wholedim}
\begin{itemize}
\item Chiyan Chen and Hongwei Xi,
Combining Programming with Theorem Proving. In Proceedings of the 10th ACM
SIGPLAN International Conference on Functional Programming (ICFP '05),
Tallinn, Estonia, September 2005.
%%%
\item Hongwei Xi,
Applied Type System, In post-workshop proceedings of TYPES 2003, 
Torino, Italy, April/May 2003.
%%%
\item Hongwei Xi, Chiyan Chen and Gang Chen,
Guarded Recursive Datatype Constructors, In Proceedings of ACM SIGPLAN
Symposium on Principles of Programming Languages (POPL '03), pp. 214-227,
New Orleans, January 2003.
%%%
\item Hongwei Xi, Dependent Types for Program Termination Verification,
In the Journal of Higher-Order and Symbolic Computation, vol. 15,
pp. 91--132, 2002.
%%%
\item Hongwei Xi and Frank Pfenning, Dependent Types in Practical
Programming. In Proceedings of ACM SIGPLAN Symposium on Principles of
Programming Languages (POPL '99), pp. 214-227, San Antonio, January 1999.
%%%

%%%
\end{itemize}\end{minipage}
} \\[\middim]
\end{tabular}

\begin{tabular}{rp{\rightdim}}
\hbox to \leftdim{\itemfont Synergistic Activities\hss} & \\[\middim]
\multicolumn{2}{c}{%
\begin{minipage}{\wholedim}
Membership in Association of Computing Machinery (ACM),
Special Interest Group in Programming Languages (SIGPLAN),
Mathematical Association of America (MAA).
\end{minipage}
} \\[\smalldim]\\[\middim]

\hbox to \leftdim{\itemfont Other Publications\hss} & \\[\middim]

\multicolumn{2}{c}{%
\begin{minipage}{\wholedim}
\begin{itemize}
%%%
\item
Dengping Zhu and Hongwei Xi,
Safe Programming with Pointers through Stateful Views. In Proceedings of
the 7th International Symposium on Practical Aspects of Declarative
Languages (PADL'05), LNCS vol. 3350, pp. 83--97, Long Beach, CA, January
2005.
\item Chiyan Chen and Hongwei Xi,
Meta-Programming through Typeful Code Representation. In Proceedings of the
8th ACM SIGPLAN International Conference on Functional Programming
(ICFP '03), pp. 275--286, Uppsala, Sweden, August 2003.
%%%
\item Hongwei Xi and Robert Harper, A Dependently Typed Assembly Language.
In Proceedings of the 6th International Conference on Functional
Programming (ICFP '01), pp. 169-180, Florence Italy, September 2001.
%%%
\item Hongwei Xi, Imperative Programming with Dependent Types,
In Proceedings of IEEE Symposium on Logic in Computer Science
(LICS '00), pp. 375-387, Santa Barbara, CA, June 2000.
%%%
%\item Hongwei Xi and Songtao Xia, Towards Array Bound Check Elimination in Java Virtual Machine Language. In Proceedings of CASCON '99, Mississauga, Ontario, November 1999.
%%%
%\item Hongwei Xi, Dependently Typed Data Structures. In Proceedings of Workshop on Algorithmic Aspects of Advanced Programming Languages (WAAAPL '99), Paris, September 1999.
%%%
%\item Femke van Raamsdonk, Paula Severi, Morten H. Sorensen and Hongwei Xi, Perpetual Reductions in lambda-calculus, Journal of Information and Computation, 149(2), pp. 173-225, March 1999.
%%%
%\item Hongwei Xi, Upper bounds for standardization and an application, Journal of Symbolic Logic, 64(1), pp. 291-303, March 1999.
%%%
%\item Hongwei Xi, Dead Code Elimination through Dependent Types. In Proceedings of the First International Workshop on Practical Aspects of Declarative Languages (PADL '99), Springer-Verlag LNCS vol. 1551, pp.228-242, San Antonio, January 1999. 
%%%
\item Hongwei Xi and Frank Pfenning, Eliminating Array Bound Checking Through Dependent Types. In Proceedings of ACM SIGPLAN Conference on Programming Language Design and Implementation (PLDI '98), pp. 249-257, Montreal, June 1998.
%%%
%\item Hongwei Xi, Towards automated termination proofs through
%"Freezing". In the proceedings of 9th International Conference on Rewriting
%Techniques and Applications (RTA '98), pp. 271-285, Japan, 1998.
%%%%
%\item Hongwei Xi, Evaluation under lambda-abstraction. In Proceedings of
%9th International Symposium on Programming Languages, Implementations,
%Logics, and Programs (PLILP '97), LNCS 1292, pp. 259-273, Southampton, 1997.
%%%%
%%\item Hongwei Xi, Simulating Eta-Expansions with Beta-Reductions in the Second-Order Polymorphic Lambda-Calculus. In the proceedings of Symposium on Logical Foundations of Computer Science (LFCS '97), LNCS, vol. 1234, pp. 399-409, Yaroslavl, 1997.
%%%%
%\item Peter B. Andrews, Matthew Bishop, Sunil Issar, Dan Nesmith, Frank Pfenning, and Hongwei Xi. TPS: A Theorem Proving System for Classical Type Theory, Journal of Automated Reasoning, 16(3), pp. 321-353, 1996.
\end{itemize}\end{minipage}
} \\[\smalldim]\\[\middim]

\hbox to \leftdim{\itemfont Advisor, Advisees and Co-authors\hss} & \\[\middim]
\multicolumn{2}{c}{%
\begin{minipage}{\wholedim}
\begin{tabular}{ll}
Advisor: & Frank Pfenning (Carnegie Mellon University) \\
Advisees (partial list): & Chiyan Chen (Boston University) \\
          & Dengping Zhu (Boston University) \\
          & Rui Shi (Boston University) \\
          & Kevin Donnelly (Boston University) \\
          & Rick Lavoie (Boston University) \\
co-authors (partial list): & Peter B. Andrews, Matthew Bishop, Chiyan Chen, \\
            & Kevin Donnelly, Robert Harper, Frank Pfenning, \\
            & Rui Shi, Walid Taha, Dengping Zhu, etc. \\
\end{tabular}
\end{minipage}
}
\end{tabular}

\end{document}
